\section{Finite lower-quota overlap below quadratic height via sparse Edgeworth}\label{sec:edgeworth-finite-lower-overlap}

Assume the uniform sparse fixed-profile Edgeworth input in \hyperref[sec:edgeworth-bounded-rank-ratio]{the covariance-transfer note}. Let $d\to\infty$, $b=d^2/n\to\infty$, $b=o(d)$, $r/b\to\infty$, $1\le r<d$, and $k/d\to\beta\in(0,\infty)$. For $a_i$ the original capped counts, put $\Gamma_i=a_i^2/(a_{i-1}a_{i+1})$, $t_k=\sqrt{a_{k-1}/a_{k+1}}$, and let $w_x$ be the positive occupied tilt with mean $dx/n$. Let $v_k$ be the first positive hitting point of phase $k\pi$ in the \hyperref[sec:analytic-angular-signs]{angular convention}. Then, for every fixed integer $j$,

\begin{equation*}
\log\Gamma_{k+j}\to1/\beta,\qquad
\log t_k=d\log w_k+\frac{d}{2k}+o(1),\qquad
d\log v_k=\log t_k+o(1).
\end{equation*}

In particular both the geometric and angular test points give the correct sign $(-1)^k$. No height assumption stronger than $n/d\to\infty$ is used within these sequences.

\subsection{Switching and scalar clipping}

At bounded offsets $i=k+O(1)$ put $m=di/n\sim\beta b$, $M=di$, and $L=\lceil B(m+\log n)\rceil=o(d)$ for sufficiently large fixed $B$. Both distances to the quota endpoints diverge since $D/d\sim r/b\to\infty$. Under uniform balanced capped arrays let $N_s$ count those with $s$ ones in row one. Move one of its ones downward in the same column to a nonfull row with no one in that column. There are at least $s(n-i-M/r)$ allowed moves because at most $i$ rows are occupied in the column and at most $M/r$ rows are full. Each array with first row size $s-1$ has at most $(d-s+1)i$ possible reverse moves. Thus $N_s\le(2m/s)N_{s-1}$ for large $d$, since $i/n=m/d=o(1)$ and $m/r=o(1)$. The factorial bound $N_s/\sum_jN_j\le(2m)^s/s!$ and a union bound over rows give $\mathbb P(\max s_j>L)=o(1)$, with an arbitrarily large fixed inverse-polynomial bound by increasing $B$.

For the scalar tuned rows, the Poisson Bernstein bound from the covariance-transfer proof, at $L\gg m+\log n$, removes any prescribed polynomial tail even with fixed polynomial weights. Conditioning on the total, which has probability at least $1/(nr+1)$, therefore leaves $o(1)$ tail probability. The clipping error in $\mathbb E[(S-m)^2]$ is $o(m)$, for example by the support bound $d$ and arbitrary polynomial tail. Applying the covariance transfer on these retained profiles yields

\begin{equation*}
a_i=A_i\sqrt d(2\pi\Lambda_*(i))^{-\nu}
  \exp[-1/(12\beta)+\nu V_i/(m(d-m))+o(1)]
\end{equation*}

where $A_i=[z^{di}](\sum_{s\le r}\binom ds z^s)^n$, $\Lambda_*(i)=i(d-m)/(d-1)$, $\nu=(d-1)/2$, and $V_i$ is the scalar size variance. This is where the Edgeworth input replaces Greenhill--McKay--Wang and removes the earlier $b=o(d^{1/3})$ restriction.

At any tilt with $m\asymp b$, the untruncated binomial mean $m_0$ belongs to $[m,2m]$ eventually: the truncated mean at untruncated mean $2m$ is $(2+o(1))m$, because $r/m\to\infty$ and the removed tail has order $o(m^{-A})$ for every fixed $A$. Monotonicity of the truncated mean gives the upper bound. Fixed raw deleted moments satisfy that tail estimate by the binomial exponential moment at a fixed positive argument, and hence $m_0=m+o(1)$, $V_i\sim m$, and central absolute moments through order four are bounded by the usual binomial scales $O(m^{j/2})$ for $2\le j\le4$. These estimates also hold uniformly over bounded real quota shifts. Therefore the variance correction in the displayed coefficient formula tends to $1/2$.

Write $E(x)=n\log H(w_x)-dx\log w_x$. The \href{https://github.com/ProgramComputer/open-math-research/blob/78e9ef2c0c0b7a5daed126efbe12a46d030d8200/topics/balanced-arrays/dense-regime/uniform-lattice-lemma.md}{integer-mean local limit} gives $A_i=e^{E(i)}(2\pi nV_i)^{-1/2}(1+o(1))$. On fixed translates $E'(x)=-d\log w_x$ and $E''(x)=-b/V_x\to-1/\beta$ uniformly. The lattice prefactor has log differences $o(1)$ since neighboring variances have ratio tending to one. The covariance prefactor has log first difference $-\nu/i+o(1)$ and second difference $o(1)$. In the signed centered first difference of $-\log a_i$, the locally constant limiting entropy curvature cancels between the two halves; this gives $\log t_k=d\log w_k+d/(2k)+o(1)$. The positive second log differences tend to $1/\beta$.

\subsection{Phase overlap at every crossing}

For any occupied tilt crossing phase $k\pi$, write $\mu=\mathbb ES$, $y=n\mu/d$, $T=nr/d$, $\varphi=\pi/d$, $G(s)=s(d-s)/[d(d-1)]$, and $a=\mathbb E G(S)$. The global phase bound $|\Psi/\pi-y|\le C d+C\min(y,T-y)/\sqrt d$ from \href{https://github.com/ProgramComputer/open-math-research/blob/4fb624ba937a8eecc1361a6066185cd81ab7b172/topics/balanced-arrays/quadratic/proof.md}{the earlier phase argument} first gives $y=O(d)$, hence $\mu=O(b)$. At this step $\mu$ need not yet be comparable to $b$. The \href{https://github.com/ProgramComputer/open-math-research/blob/78e9ef2c0c0b7a5daed126efbe12a46d030d8200/topics/balanced-arrays/dense-regime/rank-dependent/small-variance.md}{uniform third moment estimates} give $\mathbb E|S-\mu|^3\le C V\sqrt{1+V}$, $V\le\mu$, so the radial phase error multiplied by $n$ is $O(n(1+b)^{3/2}/d^3)=o(1)$ by centering at $\mu$. Even the coarser raw third-moment scale would give $O(b^2/d)=o(d)$, sufficient here. Also $a\asymp\mu/d$ and $\mathbb E[G(S)S]/a\le C(\mu+1)$, using $G(s)\le C s/d$ and $\mathbb E S^2=\mu^2+V$. The characteristic functions of both size laws are within $O((1+b)/d)=o(1)$ of one at $\varphi$ by their first moments. Thus the determinant argument is $O((1+b)/d)$ and costs $O(1+b)=o(d)$ in the phase. It follows that $y=k+o(d)$ and $\mu/b\to\beta$ at every crossing.

Now use the raw binomial-tail comparison above at $\mu\sim\beta b$. Set $Z=S-\mu$ and let $\mathbb E_G$ denote size averaging weighted by $G/a$. Lipschitz continuity of $G$ and the fourth centered moment bound give

\begin{equation*}
\mathbb E_G|Z|^3
\le \frac{G(\mu)\mathbb E|Z|^3+(C/d)\mathbb E|Z|^4}{a}
=O(b^{3/2}).
\end{equation*}

The determinant's linear phase shift is
$\mathbb E_G S-\mu=[(d-2\mu)V-\mathbb E Z^3]/[\mu(d-\mu)-V]\to1$
since $V\sim\mu$, $\mu=o(d)$ and $\mathbb E|Z|^3=O(\mu^{3/2})$. The centered third absolute moments about the means of both laws are $O(b^{3/2})$, so their linear argument expansions have cubic phase errors $O(b^{3/2}/d^3)$. Multiplied by $n$ for the radial phase, this is $O(\sqrt b/d)=o(1)$, and multiplied by $\nu$ for the determinant it is $O(b^{3/2}/d^2)=o(1)$. Thus $\Psi/\pi=y-\tfrac12+o(1)$ and $y(v_k)=k+\tfrac12+o(1)$. Integrating the identity $(d\log w_x)'=b/V_x\to1/\beta$ over this bounded shift proves $d\log v_k=\log t_k+o(1)$.

At every fixed offset $j$ the term magnitude relative to the central term at $-t_k$ tends to $e^{-j^2/(2\beta)}$, with the same limit after the $e^{o(1)}$ radius perturbation. The orientation inequality makes the ratios strictly decrease outward, and the first outward magnitude ratios stay below one. The two alternating tails are bounded by their first omitted terms. Truncate at a fixed distance, take the parameter limit and then let the truncation grow to obtain the positive Gaussian alternating sum by Poisson summation. Hence both nodes have the required sign. Furthermore, this overlap separates adjacent geometric points by $\log(t_{k+1}/t_k)\to1/\beta$, so substitutions of angular points with logarithmic discrepancy $o(1)$ preserve a join here.
